% Options for packages loaded elsewhere
\PassOptionsToPackage{unicode}{hyperref}
\PassOptionsToPackage{hyphens}{url}
\PassOptionsToPackage{dvipsnames,svgnames,x11names}{xcolor}
%
\documentclass[
  12pt,
  letterpaper,
]{book}

\usepackage{amsmath,amssymb}
\usepackage{iftex}
\ifPDFTeX
  \usepackage[T1]{fontenc}
  \usepackage[utf8]{inputenc}
  \usepackage{textcomp} % provide euro and other symbols
\else % if luatex or xetex
  \usepackage{unicode-math}
  \defaultfontfeatures{Scale=MatchLowercase}
  \defaultfontfeatures[\rmfamily]{Ligatures=TeX,Scale=1}
\fi
\usepackage{lmodern}
\ifPDFTeX\else  
    % xetex/luatex font selection
\fi
% Use upquote if available, for straight quotes in verbatim environments
\IfFileExists{upquote.sty}{\usepackage{upquote}}{}
\IfFileExists{microtype.sty}{% use microtype if available
  \usepackage[]{microtype}
  \UseMicrotypeSet[protrusion]{basicmath} % disable protrusion for tt fonts
}{}
\makeatletter
\@ifundefined{KOMAClassName}{% if non-KOMA class
  \IfFileExists{parskip.sty}{%
    \usepackage{parskip}
  }{% else
    \setlength{\parindent}{0pt}
    \setlength{\parskip}{6pt plus 2pt minus 1pt}}
}{% if KOMA class
  \KOMAoptions{parskip=half}}
\makeatother
\usepackage{xcolor}
\usepackage[margin=1in]{geometry}
\setlength{\emergencystretch}{3em} % prevent overfull lines
\setcounter{secnumdepth}{5}
% Make \paragraph and \subparagraph free-standing
\makeatletter
\ifx\paragraph\undefined\else
  \let\oldparagraph\paragraph
  \renewcommand{\paragraph}{
    \@ifstar
      \xxxParagraphStar
      \xxxParagraphNoStar
  }
  \newcommand{\xxxParagraphStar}[1]{\oldparagraph*{#1}\mbox{}}
  \newcommand{\xxxParagraphNoStar}[1]{\oldparagraph{#1}\mbox{}}
\fi
\ifx\subparagraph\undefined\else
  \let\oldsubparagraph\subparagraph
  \renewcommand{\subparagraph}{
    \@ifstar
      \xxxSubParagraphStar
      \xxxSubParagraphNoStar
  }
  \newcommand{\xxxSubParagraphStar}[1]{\oldsubparagraph*{#1}\mbox{}}
  \newcommand{\xxxSubParagraphNoStar}[1]{\oldsubparagraph{#1}\mbox{}}
\fi
\makeatother

\usepackage{color}
\usepackage{fancyvrb}
\newcommand{\VerbBar}{|}
\newcommand{\VERB}{\Verb[commandchars=\\\{\}]}
\DefineVerbatimEnvironment{Highlighting}{Verbatim}{commandchars=\\\{\}}
% Add ',fontsize=\small' for more characters per line
\usepackage{framed}
\definecolor{shadecolor}{RGB}{241,243,245}
\newenvironment{Shaded}{\begin{snugshade}}{\end{snugshade}}
\newcommand{\AlertTok}[1]{\textcolor[rgb]{0.68,0.00,0.00}{#1}}
\newcommand{\AnnotationTok}[1]{\textcolor[rgb]{0.37,0.37,0.37}{#1}}
\newcommand{\AttributeTok}[1]{\textcolor[rgb]{0.40,0.45,0.13}{#1}}
\newcommand{\BaseNTok}[1]{\textcolor[rgb]{0.68,0.00,0.00}{#1}}
\newcommand{\BuiltInTok}[1]{\textcolor[rgb]{0.00,0.23,0.31}{#1}}
\newcommand{\CharTok}[1]{\textcolor[rgb]{0.13,0.47,0.30}{#1}}
\newcommand{\CommentTok}[1]{\textcolor[rgb]{0.37,0.37,0.37}{#1}}
\newcommand{\CommentVarTok}[1]{\textcolor[rgb]{0.37,0.37,0.37}{\textit{#1}}}
\newcommand{\ConstantTok}[1]{\textcolor[rgb]{0.56,0.35,0.01}{#1}}
\newcommand{\ControlFlowTok}[1]{\textcolor[rgb]{0.00,0.23,0.31}{\textbf{#1}}}
\newcommand{\DataTypeTok}[1]{\textcolor[rgb]{0.68,0.00,0.00}{#1}}
\newcommand{\DecValTok}[1]{\textcolor[rgb]{0.68,0.00,0.00}{#1}}
\newcommand{\DocumentationTok}[1]{\textcolor[rgb]{0.37,0.37,0.37}{\textit{#1}}}
\newcommand{\ErrorTok}[1]{\textcolor[rgb]{0.68,0.00,0.00}{#1}}
\newcommand{\ExtensionTok}[1]{\textcolor[rgb]{0.00,0.23,0.31}{#1}}
\newcommand{\FloatTok}[1]{\textcolor[rgb]{0.68,0.00,0.00}{#1}}
\newcommand{\FunctionTok}[1]{\textcolor[rgb]{0.28,0.35,0.67}{#1}}
\newcommand{\ImportTok}[1]{\textcolor[rgb]{0.00,0.46,0.62}{#1}}
\newcommand{\InformationTok}[1]{\textcolor[rgb]{0.37,0.37,0.37}{#1}}
\newcommand{\KeywordTok}[1]{\textcolor[rgb]{0.00,0.23,0.31}{\textbf{#1}}}
\newcommand{\NormalTok}[1]{\textcolor[rgb]{0.00,0.23,0.31}{#1}}
\newcommand{\OperatorTok}[1]{\textcolor[rgb]{0.37,0.37,0.37}{#1}}
\newcommand{\OtherTok}[1]{\textcolor[rgb]{0.00,0.23,0.31}{#1}}
\newcommand{\PreprocessorTok}[1]{\textcolor[rgb]{0.68,0.00,0.00}{#1}}
\newcommand{\RegionMarkerTok}[1]{\textcolor[rgb]{0.00,0.23,0.31}{#1}}
\newcommand{\SpecialCharTok}[1]{\textcolor[rgb]{0.37,0.37,0.37}{#1}}
\newcommand{\SpecialStringTok}[1]{\textcolor[rgb]{0.13,0.47,0.30}{#1}}
\newcommand{\StringTok}[1]{\textcolor[rgb]{0.13,0.47,0.30}{#1}}
\newcommand{\VariableTok}[1]{\textcolor[rgb]{0.07,0.07,0.07}{#1}}
\newcommand{\VerbatimStringTok}[1]{\textcolor[rgb]{0.13,0.47,0.30}{#1}}
\newcommand{\WarningTok}[1]{\textcolor[rgb]{0.37,0.37,0.37}{\textit{#1}}}

\providecommand{\tightlist}{%
  \setlength{\itemsep}{0pt}\setlength{\parskip}{0pt}}\usepackage{longtable,booktabs,array}
\usepackage{calc} % for calculating minipage widths
% Correct order of tables after \paragraph or \subparagraph
\usepackage{etoolbox}
\makeatletter
\patchcmd\longtable{\par}{\if@noskipsec\mbox{}\fi\par}{}{}
\makeatother
% Allow footnotes in longtable head/foot
\IfFileExists{footnotehyper.sty}{\usepackage{footnotehyper}}{\usepackage{footnote}}
\makesavenoteenv{longtable}
\usepackage{graphicx}
\makeatletter
\def\maxwidth{\ifdim\Gin@nat@width>\linewidth\linewidth\else\Gin@nat@width\fi}
\def\maxheight{\ifdim\Gin@nat@height>\textheight\textheight\else\Gin@nat@height\fi}
\makeatother
% Scale images if necessary, so that they will not overflow the page
% margins by default, and it is still possible to overwrite the defaults
% using explicit options in \includegraphics[width, height, ...]{}
\setkeys{Gin}{width=\maxwidth,height=\maxheight,keepaspectratio}
% Set default figure placement to htbp
\makeatletter
\def\fps@figure{htbp}
\makeatother

% header.tex
\usepackage{xcolor}        % Para colores
\usepackage{polyglossia}   % Soporte de idiomas con XeLaTeX
\setdefaultlanguage{spanish}

% Definición de colores personalizados
\definecolor{fecundidad}{HTML}{FFC125}      % Amarillo
\definecolor{supervivencia}{HTML}{36648B}   % Azul
\definecolor{crecimiento}{HTML}{548B54}     % Verde
\definecolor{retroceso}{HTML}{7A378B}       % Violeta
\definecolor{permanencia}{HTML}{26466A}     % Azul oscuro
\usepackage{booktabs}
\usepackage{longtable}
\usepackage{array}
\usepackage{multirow}
\usepackage{wrapfig}
\usepackage{float}
\usepackage{colortbl}
\usepackage{pdflscape}
\usepackage{tabu}
\usepackage{threeparttable}
\usepackage{threeparttablex}
\usepackage[normalem]{ulem}
\usepackage{makecell}
\usepackage{xcolor}
\usepackage{caption}
\usepackage{anyfontsize}
\makeatletter
\@ifpackageloaded{caption}{}{\usepackage{caption}}
\AtBeginDocument{%
\ifdefined\contentsname
  \renewcommand*\contentsname{Table of contents}
\else
  \newcommand\contentsname{Table of contents}
\fi
\ifdefined\listfigurename
  \renewcommand*\listfigurename{List of Figures}
\else
  \newcommand\listfigurename{List of Figures}
\fi
\ifdefined\listtablename
  \renewcommand*\listtablename{List of Tables}
\else
  \newcommand\listtablename{List of Tables}
\fi
\ifdefined\figurename
  \renewcommand*\figurename{Figure}
\else
  \newcommand\figurename{Figure}
\fi
\ifdefined\tablename
  \renewcommand*\tablename{Table}
\else
  \newcommand\tablename{Table}
\fi
}
\@ifpackageloaded{float}{}{\usepackage{float}}
\floatstyle{ruled}
\@ifundefined{c@chapter}{\newfloat{codelisting}{h}{lop}}{\newfloat{codelisting}{h}{lop}[chapter]}
\floatname{codelisting}{Listing}
\newcommand*\listoflistings{\listof{codelisting}{List of Listings}}
\makeatother
\makeatletter
\makeatother
\makeatletter
\@ifpackageloaded{caption}{}{\usepackage{caption}}
\@ifpackageloaded{subcaption}{}{\usepackage{subcaption}}
\makeatother

\ifLuaTeX
  \usepackage{selnolig}  % disable illegal ligatures
\fi
\usepackage{bookmark}

\IfFileExists{xurl.sty}{\usepackage{xurl}}{} % add URL line breaks if available
\urlstyle{same} % disable monospaced font for URLs
\hypersetup{
  colorlinks=true,
  linkcolor={Maroon},
  filecolor={Maroon},
  citecolor={Blue},
  urlcolor={Blue},
  pdfcreator={LaTeX via pandoc}}


\author{}
\date{}

\begin{document}
\frontmatter

\renewcommand*\contentsname{Table of contents}
{
\hypersetup{linkcolor=}
\setcounter{tocdepth}{2}
\tableofcontents
}

\mainmatter
\chapter{COMPADRE y Orquídeas}\label{Compadre}

Por: Raymond L. Tremblay

\subsubsection{Librerías de R requeridas para el siguiente
módulo}\label{libreruxedas-de-r-requeridas-para-el-siguiente-muxf3dulo}

\begin{Shaded}
\begin{Highlighting}[]
\FunctionTok{library}\NormalTok{(Rcompadre) }\CommentTok{\# Paquete para trabajar con la base de datos de COMPADRE y COMADRE}
\FunctionTok{library}\NormalTok{(tidyverse) }\CommentTok{\# Paquete para activar multiples paquetes}
\FunctionTok{library}\NormalTok{(gt) }\CommentTok{\# Paquete para formatear tablas}
\FunctionTok{library}\NormalTok{(kableExtra) }\CommentTok{\# Paquete para formatear tablas}

\FunctionTok{library}\NormalTok{(broom)}
\FunctionTok{library}\NormalTok{(ggdist)}
\FunctionTok{library}\NormalTok{(distributional)}
\FunctionTok{library}\NormalTok{(leaflet)}
\FunctionTok{library}\NormalTok{(Rage) }\CommentTok{\# función para calcular el largo de vida esperada entre otras}
\FunctionTok{library}\NormalTok{(car)}
\FunctionTok{library}\NormalTok{(popdemo)}
\FunctionTok{library}\NormalTok{(WRS2)}
\end{Highlighting}
\end{Shaded}

Uno de los esfuerzos más grandes para compilar el conocimiento actual
sobre la ecología de poblaciones es la creación y mantenimiento de las
bases de datos de COMPADRE y COMADRE. COMPADRE es una base de datos de
matrices de transición de estados de poblaciones de plantas y COMADRE es
una base de datos de matrices de transición de estados de poblaciones de
animales. Estas bases de datos son una fuente de información para
evaluar la dinámica de poblaciones de plantas y animales. En este
capítulo se presenta una introducción a la base de datos de COMPADRE y
algunas de las funciones en los paquetes de \textbf{Rcompadre}. En el
próximo capítulo se presentarán funciones incluidas en el paquete
\textbf{Rage}. La información aquí no sustituye la información provista
en la página web de COMPADRE
\textless{}\url{https://compadre-db.org/}\textgreater{} y los tutoriales
provistos en la sección de Education de la misma
página\textless{}\url{https://compadre-db.org/Education}\textgreater.

El valor de las bases de datos reside en tener en un mismo sitio una
gran cantidad de información ~organizada de manera uniforme y accesible,
además de que su formato permite utilizarla directamente en diferentes
análisis. En el caso de COMPADRE y COMADRE, la información está
organizada en matrices de transición de estados de poblaciones (MPM) que
se pueden usar para evaluar la dinámica de poblaciones de plantas y
animales y con múltiples variables, tal como sus coordenadas
geográficas, el tipo de hábitat, la temporalidad de los datos, la fuente
de la información, entre otros.

En este capítulo se aprenderá a utilizar la base de datos de COMPADRE
para evaluar la dinámica de poblaciones de orquídeas. Se evaluará cómo
extraer diferentes tipos de datos para responder a preguntas
específicas.

\begin{center}\rule{0.5\linewidth}{0.5pt}\end{center}

\subsection{Acceso a los datos de
COMPADRE.}\label{acceso-a-los-datos-de-compadre.}

\subsubsection{Descargando el archivo de
datos}\label{descargando-el-archivo-de-datos}

Hay dos métodos de tener acceso a los datos de COMPADRE (datos de
dinámica poblacional de plantas) y COMADRE (datos de dinámica
poblacional de animales). Se puede acceder a los datos directamente de
la página web o bajar el archivo de datos desde el siguiente enlace
\textless{}\url{https://compadre-db.org/Data/Compadre}\textgreater{}
seleccionando la pestaña de \textbf{Download COMPADRE} o
\textbf{Download COMADRE}.

Para tener acceso a los datos cuando no cuente con conexión a internet
es necesario descargar el archivo en su ordenador y cargarlos en su
proyecto de RStudio. Al momento de la edición de este libro la versión
más reciente de COMPADRE es la versión 6.28.8.0 (fecha de la base de
datos: agosto\_14\_2025) y se puede descargar en el siguiente enlace:
\textless\textless{}\url{https://compadre-db.org/Data/Compadre/COMPADRE_v}.

\paragraph{Acceso a los datos desde su computadora
(ordenador)}\label{acceso-a-los-datos-desde-su-computadora-ordenador}

\begin{itemize}
\tightlist
\item
  USAR el código siguiente si ha descargado los datos en su computadora.
\end{itemize}

load(dirección/al/archivo/a/su/proyecto/de/RStudio) si tiene el
repositorio en su computadora (preferiblemente en su proyecto de
RStudio)

\begin{center}\rule{0.5\linewidth}{0.5pt}\end{center}

\subsubsection{Carga de datos con el enlace de
web}\label{carga-de-datos-con-el-enlace-de-web}

\begin{itemize}
\tightlist
\item
  Usando la función \emph{cdb\_fetch()} se pueden descargar los datos
  directamente de la página web de COMPADRE. Esta función se encuentra
  en el paquete de \textbf{Rcompadre}.
\end{itemize}

\begin{Shaded}
\begin{Highlighting}[]
\NormalTok{compadre }\OtherTok{\textless{}{-}} \FunctionTok{cdb\_fetch}\NormalTok{(}\StringTok{"compadre"}\NormalTok{) }\CommentTok{\# Usar este código para bajar los datos del repositorio de COMPADRE.}
\end{Highlighting}
\end{Shaded}

\begin{center}\rule{0.5\linewidth}{0.5pt}\end{center}

\subsubsection{Renombrar la base datos a algo
sencillo}\label{renombrar-la-base-datos-a-algo-sencillo}

Se considera una buena práctica renombrar los objetos de forma tal que
sea simple llamarlos posteriormente en nuestro script. Para ello, en
este ejemplo se renombrará la base de datos de COMPADRE a algo más
sencillo \textbf{TodasSp}, reemplazando \textbf{compadre}. - Se
visualiza el nombre de cada columna en la base de datos. - Hay 59
variables en la base de datos de COMPADRE.

\begin{Shaded}
\begin{Highlighting}[]
\NormalTok{TodasSp}\OtherTok{=}\NormalTok{compadre}
\CommentTok{\#names(TodasSp) \# La lista de los nombres de las variables en el archivo COMPADRE, Nota que 59 variables en la base de datos de COMPADRE.}
\end{Highlighting}
\end{Shaded}

\begin{itemize}
\tightlist
\item
  Conocer la dimensión de la base de dato. Al momento de la edición de
  este libro la base de datos de COMPADRE tiene 9149 filas y 59 columnas
  (variables).
\end{itemize}

\begin{Shaded}
\begin{Highlighting}[]
\FunctionTok{dim}\NormalTok{(TodasSp) }
\end{Highlighting}
\end{Shaded}

\begin{verbatim}
[1] 9146   59
\end{verbatim}

\begin{center}\rule{0.5\linewidth}{0.5pt}\end{center}

\subsection{Explorando la base de datos de
COMPADRE}\label{explorando-la-base-de-datos-de-compadre}

Se estará evaluando múltiples funciones en Rcompadre para aprender a
usar la base de datos de COMPADRE.

Para evaluar si la base de datos contiene información sobre especies
particulares de orquídeas, ~se utilizará la función
`cdb\_check\_species'. Es relevante mencionar que se puede evaluar
múltiples especies de forma simultánea. Aquí se evaluó si la base
contiene información sobre las especies \emph{Epidendrum xanthinum},
\emph{Epipactis gigantea} y \emph{Caladenia latifolia}.

Como podrán notar, solamente existen datos para una de las tres que
probamos en la base de datos de COMPADRE.

\begin{Shaded}
\begin{Highlighting}[]
\NormalTok{especies}\OtherTok{\textless{}{-}} \FunctionTok{c}\NormalTok{(}\StringTok{"Epidendrum xanthinum"}\NormalTok{, }\StringTok{"Epipactis gigantea"}\NormalTok{, }\StringTok{"Caladenia latifolia"}\NormalTok{)}
\FunctionTok{cdb\_check\_species}\NormalTok{(TodasSp, especies) }\SpecialCharTok{|\textgreater{}} \FunctionTok{gt}\NormalTok{()}
\end{Highlighting}
\end{Shaded}

\begin{table}
\fontsize{12.0pt}{14.0pt}\selectfont
\begin{tabular*}{\linewidth}{@{\extracolsep{\fill}}lc}
\toprule
species & in\_db \\ 
\midrule\addlinespace[2.5pt]
Epidendrum xanthinum & TRUE \\ 
Epipactis gigantea & FALSE \\ 
Caladenia latifolia & FALSE \\ 
\bottomrule
\end{tabular*}
\end{table}

\begin{center}\rule{0.5\linewidth}{0.5pt}\end{center}

\subsection{Crear un subconjunto de la base de datos de
COMPADRE}\label{crear-un-subconjunto-de-la-base-de-datos-de-compadre}

Antes de seguir vamos a extraer de la base de datos de COMPADRE las
especies de orquídeas y crear un nuevo objeto con todas las especies de
orquídeas. La ventaja de crear un conjunto de datos es reducir el tamaño
del archivo y facilitar los análisis posteriores.

Seleccionaremos todas las matrices que pertenecen a la familia
Orchidaceae. Se observa que la base 788 filas, por consecuencia 788
matrices de transición de estados de poblaciones de orquídeas.

\begin{Shaded}
\begin{Highlighting}[]
\NormalTok{index\_O }\OtherTok{=}\NormalTok{ TodasSp }\SpecialCharTok{\%\textgreater{}\%}
  \FunctionTok{filter}\NormalTok{(Family }\SpecialCharTok{\%in\%} \FunctionTok{c}\NormalTok{(}\StringTok{"Orchidaceae"}\NormalTok{)) }\CommentTok{\# Extraer la orquídeas de la base de datos.}
\end{Highlighting}
\end{Shaded}

\begin{center}\rule{0.5\linewidth}{0.5pt}\end{center}

\subsection{La metadata de la base de datos de
COMPADRE}\label{la-metadata-de-la-base-de-datos-de-compadre}

Con la función \textbf{cdb\_metadata()} se puede evaluar la metadata de
la base de datos. La función muestra el número de filas y columnas, que
incluye mútliples variables, tal como el nombre de la especies usada en
la fuente de información `SpeciesAuthor', el nombre del taxón reconocido
``SpeciesAccepted'', la fuente de información como el DOI\_ISBN, los
años de estudio, sus coordenadas geográficas, el tipo de hábitat, la
fuente de la información, entre otros.

\begin{Shaded}
\begin{Highlighting}[]
\FunctionTok{cdb\_metadata}\NormalTok{(index\_O) }\SpecialCharTok{\%\textgreater{}\%} \FunctionTok{head}\NormalTok{() }\SpecialCharTok{|\textgreater{}} \FunctionTok{gt}\NormalTok{() }\CommentTok{\# Evaluar la metadata de la base de datos de COMPADRE}
\end{Highlighting}
\end{Shaded}

\begin{table}
\fontsize{12.0pt}{14.0pt}\selectfont
\begin{tabular*}{\linewidth}{@{\extracolsep{\fill}}rlllllllllllllllllllrllrrrrllllrrrrllllllllrlrrlrllllrrlll}
\toprule
MatrixID & SpeciesAuthor & SpeciesAccepted & CommonName & Kingdom & Phylum & Class & Order & Family & Genus & Species & Infraspecies & InfraspeciesType & OrganismType & DicotMonoc & AngioGymno & Authors & Journal & SourceType & OtherType & YearPublication & DOI\_ISBN & AdditionalSource & StudyDuration & StudyStart & StudyEnd & ProjectionInterval & MatrixCriteriaSize & MatrixCriteriaOntogeny & MatrixCriteriaAge & MatrixPopulation & NumberPopulations & Lat & Lon & Altitude & Country & Continent & Ecoregion & StudiedSex & MatrixComposite & MatrixSeasonal & MatrixTreatment & MatrixCaptivity & MatrixStartYear & MatrixStartSeason & MatrixStartMonth & MatrixEndYear & MatrixEndSeason & MatrixEndMonth & CensusType & MatrixSplit & MatrixFec & Observations & MatrixDimension & SurvivalIssue & \_Database & \_PopulationStatus & \_PublicationStatus \\ 
\midrule\addlinespace[2.5pt]
238285 & Caladenia\_amonea & Caladenia amonea & NA & Plantae & Tracheophyta & Liliopsida & Asparagales & Orchidaceae & Caladenia & amonea & NA & NA & Herbaceous perennial & Monocot & Angiosperm & Raymond L. Tremblay; P\\'erez; Larcombe; Brown; Quarmby; Bickerton; French; Bould & Aust J Bot & Journal Article & NA & 2009 & 10.1071/BT08167 & NA & 12 & 1996 & 2007 & 1 & No & Yes & No & North-eastern Melbourne, Vic & 1 & NA & NA & NA & AUS & Oceania & TBM & A & Mean & No & Unmanipulated & W & 1996 & Summer & 8 & 2007 & Autumn & 9 & NA & Divided & NA & Locations of surveyed areas are kept vague to reduce the likelihood of unauthorised disturbance of the populations & 3 & 1 & COMPADRE & Released & Released \\ 
238286 & Caladenia\_argocalla & Caladenia argocalla & NA & Plantae & Tracheophyta & Liliopsida & Asparagales & Orchidaceae & Caladenia & argocalla & NA & NA & Herbaceous perennial & Monocot & Angiosperm & Raymond L. Tremblay; P\\'erez; Larcombe; Brown; Quarmby; Bickerton; French; Bould & Aust J Bot & Journal Article & NA & 2009 & 10.1071/BT08167 & NA & 12 & 1996 & 2007 & 1 & No & Yes & No & Adelaide Hills and Clare Valley, SA & 1 & NA & NA & NA & AUS & Oceania & TBM & A & Mean & No & Unmanipulated & W & 2003 & Autumn & 9 & 2007 & Autumn & 10 & NA & Divided & No & Locations of surveyed areas are kept vague to reduce the likelihood of unauthorised disturbance of the populations & 3 & 1 & COMPADRE & Released & Released \\ 
238287 & Caladenia\_clavigera & Caladenia clavigera & NA & Plantae & Tracheophyta & Liliopsida & Asparagales & Orchidaceae & Caladenia & clavigera & NA & NA & Herbaceous perennial & Monocot & Angiosperm & Raymond L. Tremblay; P\\'erez; Larcombe; Brown; Quarmby; Bickerton; French; Bould & Aust J Bot & Journal Article & NA & 2009 & 10.1071/BT08167 & NA & 12 & 1996 & 2007 & 1 & No & Yes & No & North-eastern Melbourne, Vic & 1 & NA & NA & NA & AUS & Oceania & TBM & A & Mean & No & Unmanipulated & W & 1997 & Autumn & 9 & 2007 & Autumn & 10 & NA & Divided & NA & Locations of surveyed areas are kept vague to reduce the likelihood of unauthorised disturbance of the populations & 3 & 1 & COMPADRE & Released & Released \\ 
238288 & Caladenia\_elegans & Caladenia elegans & NA & Plantae & Tracheophyta & Liliopsida & Asparagales & Orchidaceae & Caladenia & elegans & NA & NA & Herbaceous perennial & Monocot & Angiosperm & Raymond L. Tremblay; P\\'erez; Larcombe; Brown; Quarmby; Bickerton; French; Bould & Aust J Bot & Journal Article & NA & 2009 & 10.1071/BT08167 & NA & 12 & 1996 & 2007 & 1 & No & Yes & No & Western Australia & 1 & NA & NA & NA & AUS & Oceania & TBM & A & Mean & No & Unmanipulated & W & 1998 & Summer & 7 & 2007 & Summer & 8 & NA & Divided & NA & Locations of surveyed areas are kept vague to reduce the likelihood of unauthorised disturbance of the populations & 3 & 1 & COMPADRE & Released & Released \\ 
238289 & Caladenia\_graniticola & Caladenia graniticola & NA & Plantae & Tracheophyta & Liliopsida & Asparagales & Orchidaceae & Caladenia & graniticola & NA & NA & Herbaceous perennial & Monocot & Angiosperm & Raymond L. Tremblay; P\\'erez; Larcombe; Brown; Quarmby; Bickerton; French; Bould & Aust J Bot & Journal Article & NA & 2009 & 10.1071/BT08167 & NA & 12 & 1996 & 2007 & 1 & No & Yes & No & Western Australia & 1 & NA & NA & NA & AUS & Oceania & TBM & A & Mean & No & Unmanipulated & W & 2004 & Spring & 9 & 2007 & Spring & 9 & NA & Divided & NA & Locations of surveyed areas are kept vague to reduce the likelihood of unauthorised disturbance of the populations & 3 & 1 & COMPADRE & Released & Released \\ 
238290 & Caladenia\_macroclavia & Caladenia macroclavia & NA & Plantae & Tracheophyta & Liliopsida & Asparagales & Orchidaceae & Caladenia & macroclavia & NA & NA & Herbaceous perennial & Monocot & Angiosperm & Raymond L. Tremblay; P\\'erez; Larcombe; Brown; Quarmby; Bickerton; French; Bould & Aust J Bot & Journal Article & NA & 2009 & 10.1071/BT08167 & NA & 12 & 1996 & 2007 & 1 & No & Yes & No & Yorke Peninsula, SA & 1 & NA & NA & NA & AUS & Oceania & TBM & A & Mean & No & Unmanipulated & W & 2001 & Autumn & 9 & 2007 & Autumn & 10 & NA & Divided & NA & Locations of surveyed areas are kept vague to reduce the likelihood of unauthorised disturbance of the populations & 3 & 1 & COMPADRE & Released & Released \\ 
\bottomrule
\end{tabular*}
\end{table}

\begin{center}\rule{0.5\linewidth}{0.5pt}\end{center}

\subsection{Extraer matrices}\label{extraer-matrices}

¿Cómo extraer matrices de transición de estados de las poblaciones de
orquídeas?

Vamos a seleccionar solamente unas matrices de transición de estados de
poblaciones de orquídeas. En este caso se selecciona la primera fila del
objeto \textbf{index\_O} que corresponde a la especie
\textbf{\emph{Laelia speciosa}}.

Observamos que hay 4 matriz en la lista. La matriz de transición de
estados de poblaciones (matU), la matriz de fecundidad (matF), la matriz
de clonación (matC) y la suma de las tres matrices
(matA=matU+MatF+matC).

\begin{Shaded}
\begin{Highlighting}[]
\NormalTok{index\_O }\SpecialCharTok{\%\textgreater{}\%} \FunctionTok{filter}\NormalTok{(SpeciesAccepted }\SpecialCharTok{\%in\%} \FunctionTok{c}\NormalTok{(}\StringTok{"Laelia speciosa"}\NormalTok{)) }\OtherTok{{-}\textgreater{}}\NormalTok{ Laeliasp }\CommentTok{\# Extraer la información de la especie }

\FunctionTok{head}\NormalTok{(}\FunctionTok{cdb\_metadata}\NormalTok{(Laeliasp), }\AttributeTok{n =} \DecValTok{3}\NormalTok{) }\SpecialCharTok{|\textgreater{}} 
  \FunctionTok{gt}\NormalTok{()  }\CommentTok{\# Visualizar la información de la especie Laelia speciosa}
\end{Highlighting}
\end{Shaded}

\begin{table}
\fontsize{12.0pt}{14.0pt}\selectfont
\begin{tabular*}{\linewidth}{@{\extracolsep{\fill}}rlllllllllllllllllllrllrrrrllllrrrrllllllllrlrrlrllllrrlll}
\toprule
MatrixID & SpeciesAuthor & SpeciesAccepted & CommonName & Kingdom & Phylum & Class & Order & Family & Genus & Species & Infraspecies & InfraspeciesType & OrganismType & DicotMonoc & AngioGymno & Authors & Journal & SourceType & OtherType & YearPublication & DOI\_ISBN & AdditionalSource & StudyDuration & StudyStart & StudyEnd & ProjectionInterval & MatrixCriteriaSize & MatrixCriteriaOntogeny & MatrixCriteriaAge & MatrixPopulation & NumberPopulations & Lat & Lon & Altitude & Country & Continent & Ecoregion & StudiedSex & MatrixComposite & MatrixSeasonal & MatrixTreatment & MatrixCaptivity & MatrixStartYear & MatrixStartSeason & MatrixStartMonth & MatrixEndYear & MatrixEndSeason & MatrixEndMonth & CensusType & MatrixSplit & MatrixFec & Observations & MatrixDimension & SurvivalIssue & \_Database & \_PopulationStatus & \_PublicationStatus \\ 
\midrule\addlinespace[2.5pt]
252937 & Laelia\_speciosa\_ & Laelia speciosa & Flor de Corpus, Flor de mayo & Plantae & Tracheophyta & Liliopsida & Asparagales & Orchidaceae & Laelia & speciosa & NA & variety & Epiphyte & Eudicot & Angiosperm & Mariana Hern\\'andez-Apolinar & NA & Thesis & NA & 1992 & NA & NA & 3 & 1987 & 1989 & 1 & Yes & Yes & Yes & Cerro El Olvido & 2 & 19.43333 & -101.4667 & 2100 & MEX & S America & TBM & H & Individual & NA & Unmanipulated & W & 1987 & Spring & 5 & 1989 & Winter & 1 & Birth-pulse & Indivisible & Yes & This site was disturbe by ilegal extraction, both matrices were presented & 5 & NA & COMPADRE & Released & Released \\ 
252938 & Laelia\_speciosa\_ & Laelia speciosa & Flor de Corpus, Flor de mayo & Plantae & Tracheophyta & Liliopsida & Asparagales & Orchidaceae & Laelia & speciosa & NA & variety & Epiphyte & Eudicot & Angiosperm & Mariana Hern\\'andez-Apolinar & NA & Thesis & NA & 1992 & NA & NA & 3 & 1987 & 1989 & 1 & Yes & Yes & Yes & Cerro El Olvido & 2 & 19.43333 & -101.4667 & 2100 & MEX & S America & TBM & H & Individual & NA & Unmanipulated, anthropic & W & 1987 & Spring & 5 & 1989 & Winter & 1 & Birth-pulse & Indivisible & Yes & Matrix includes data after ilegal extraction of whole plants and/or the two most recent pseudobulbs, including the newest in bloom & 5 & NA & COMPADRE & Released & Released \\ 
252939 & Laelia\_speciosa\_ & Laelia speciosa & Flor de Corpus, Flor de mayo & Plantae & Tracheophyta & Liliopsida & Asparagales & Orchidaceae & Laelia & speciosa & NA & variety & Epiphyte & Eudicot & Angiosperm & Mariana Hern\\'andez-Apolinar & NA & Thesis & NA & 1992 & NA & NA & 3 & 1987 & 1989 & 1 & Yes & Yes & Yes & Cerro El Olvido & 2 & 19.43333 & -101.4667 & 2100 & MEX & S America & TBM & H & Individual & NA & Unmanipulated & W & 1987 & Spring & 5 & 1989 & Winter & 1 & Birth-pulse & Indivisible & Yes & Matrix of Conserved site that includes only 4 stage categories. Seed category was not included. & 4 & NA & COMPADRE & Released & Released \\ 
\bottomrule
\end{tabular*}
\end{table}

\begin{center}\rule{0.5\linewidth}{0.5pt}\end{center}

\subsection{Extraer la matriz de transición de estados de
poblaciones}\label{extraer-la-matriz-de-transiciuxf3n-de-estados-de-poblaciones}

En el siguiente script se muestra como extraer la matriz de transición
de estados de poblaciones (matA) de \textbf{Laelia speciosa} y las
etapas de esta matrices. Nota que el objeto \textbf{mat} es una lista de
matrices que incluye la matriz de transición (matU) de estados de una
población, la matriz de fecundidad (matF), la matriz de clonaje (matC) y
la suma de las tres matrices (matA=matU+MatF+matC) vea el capítulo de
{[}Matrices de transición, fecundidad y clonal{]}

\begin{Shaded}
\begin{Highlighting}[]
\NormalTok{Laeliaspmat }\OtherTok{=} \FunctionTok{matA}\NormalTok{(Laeliasp) }\CommentTok{\# Extraer la matriz de transiciones de estados de la población (matA) de las cuatro filas}
\NormalTok{Laeliaspmat}
\end{Highlighting}
\end{Shaded}

\begin{verbatim}
[[1]]
        A1   A2   A3       A4        A5
A1 0.0e+00 0.00 0.00 52686.00 127342.00
A2 2.2e-05 0.00 0.00     0.00      0.00
A3 0.0e+00 0.71 0.93     0.00      0.00
A4 0.0e+00 0.00 0.07     0.76      0.00
A5 0.0e+00 0.00 0.00     0.24      0.75

[[2]]
        A1   A2   A3       A4        A5
A1 0.0e+00 0.00 0.00 52686.00 127324.00
A2 2.2e-05 0.00 0.00     0.00      0.00
A3 0.0e+00 0.71 0.64     0.00      0.00
A4 0.0e+00 0.00 0.05     0.61      0.00
A5 0.0e+00 0.00 0.00     0.19      0.33

[[3]]
     A1   A2   A3    A4
A1 0.00 0.00 1.64 11.25
A2 0.71 0.93 0.00  0.00
A3 0.00 0.07 0.76  0.00
A4 0.00 0.00 0.24  0.75

[[4]]
     A1   A2   A3    A4
A1 0.00 0.00 1.64 11.25
A2 0.71 0.64 0.00  0.00
A3 0.00 0.05 0.61  0.00
A4 0.00 0.00 0.19  0.33
\end{verbatim}

\begin{center}\rule{0.5\linewidth}{0.5pt}\end{center}

\section{Extraer matA, matU, matF y
matC}\label{extraer-mata-matu-matf-y-matc}

Note que esta función es muy práctica porque extrae la matriz de
transición de estados de poblaciones (matA), la matriz de supervivencia
(matU), la matriz de fecundidad (matF) y la matriz de clonaje (matC) sin
necesidad de usar el código largo en adición de los nombres de las
etapas.

Se selecciona la segunda fila del objeto \textbf{Laeliasp} que
corresponde a la especie \textbf{Laelia speciosa} y la segunda matriz en
la fila{[}{[}2{]}{]}.

\begin{Shaded}
\begin{Highlighting}[]
\NormalTok{Laeliasp}\SpecialCharTok{$}\NormalTok{mat[[}\DecValTok{2}\NormalTok{]]}
\end{Highlighting}
\end{Shaded}

\begin{verbatim}
  MatrixClassOrganized                                 MatrixClassAuthor
1               active                                          S(Seeds)
2               active Pl(Seedlings): one small pseudobulb of 3mm height
3               active                    I(Juveniles): 2-10 pseudobulbs
4               active                    II(Adult 1): 11-20 pseudobulbs
5               active                     III(Adult 2): >21 pseudobulbs

matA:
        1    2    3        4         5
1 0.0e+00 0.00 0.00 52686.00 127324.00
2 2.2e-05 0.00 0.00     0.00      0.00
3 0.0e+00 0.71 0.64     0.00      0.00
4 0.0e+00 0.00 0.05     0.61      0.00
5 0.0e+00 0.00 0.00     0.19      0.33

matU:
  1 2 3 4 5
1 0 0 0 0 0
2 0 0 0 0 0
3 0 0 0 0 0
4 0 0 0 0 0
5 0 0 0 0 0

matF:
  1 2 3 4 5
1 0 0 0 0 0
2 0 0 0 0 0
3 0 0 0 0 0
4 0 0 0 0 0
5 0 0 0 0 0

matC:
  1 2 3 4 5
1 0 0 0 0 0
2 0 0 0 0 0
3 0 0 0 0 0
4 0 0 0 0 0
5 0 0 0 0 0
\end{verbatim}

\begin{center}\rule{0.5\linewidth}{0.5pt}\end{center}

\subsubsection{Evaluar si un genero o especie este en la base de datos
de
COMPADRE}\label{evaluar-si-un-genero-o-especie-este-en-la-base-de-datos-de-compadre}

\begin{itemize}
\tightlist
\item
  Para filtrar un genero específico se puede usar el siguiente código.

  \begin{itemize}
  \tightlist
  \item
    Aquí seleccionamos todas las matrices del género \textbf{Epipactis}.
  \item
    Nota que hay 59 filas, que corresponde típicamente a 59 matrices de
    diferentes poblaciones, tiempo o un análisis teóricos.
  \end{itemize}
\end{itemize}

\begin{Shaded}
\begin{Highlighting}[]
\NormalTok{Epipactis }\OtherTok{=}\NormalTok{ TodasSp }\SpecialCharTok{\%\textgreater{}\%}
  \FunctionTok{filter}\NormalTok{(Genus }\SpecialCharTok{\%in\%} \FunctionTok{c}\NormalTok{(}\StringTok{"Epipactis"}\NormalTok{)) }\CommentTok{\# Filtrar la información de *Epipactis* en la base de datos.}

\FunctionTok{head}\NormalTok{(}\FunctionTok{cdb\_metadata}\NormalTok{(Epipactis), }\AttributeTok{n =} \DecValTok{3}\NormalTok{) }\SpecialCharTok{|\textgreater{}} 
  \FunctionTok{gt}\NormalTok{()}
\end{Highlighting}
\end{Shaded}

\begin{table}
\fontsize{12.0pt}{14.0pt}\selectfont
\begin{tabular*}{\linewidth}{@{\extracolsep{\fill}}rlllllllllllllllllllrllrrrrllllrrrrllllllllrllrllllllrrlll}
\toprule
MatrixID & SpeciesAuthor & SpeciesAccepted & CommonName & Kingdom & Phylum & Class & Order & Family & Genus & Species & Infraspecies & InfraspeciesType & OrganismType & DicotMonoc & AngioGymno & Authors & Journal & SourceType & OtherType & YearPublication & DOI\_ISBN & AdditionalSource & StudyDuration & StudyStart & StudyEnd & ProjectionInterval & MatrixCriteriaSize & MatrixCriteriaOntogeny & MatrixCriteriaAge & MatrixPopulation & NumberPopulations & Lat & Lon & Altitude & Country & Continent & Ecoregion & StudiedSex & MatrixComposite & MatrixSeasonal & MatrixTreatment & MatrixCaptivity & MatrixStartYear & MatrixStartSeason & MatrixStartMonth & MatrixEndYear & MatrixEndSeason & MatrixEndMonth & CensusType & MatrixSplit & MatrixFec & Observations & MatrixDimension & SurvivalIssue & \_Database & \_PopulationStatus & \_PublicationStatus \\ 
\midrule\addlinespace[2.5pt]
238565 & Epipactis\_atrorubens\_3 & Epipactis atrorubens & Darkred helleborine & Plantae & Tracheophyta & Liliopsida & Asparagales & Orchidaceae & Epipactis & atrorubens & NA & NA & Herbaceous perennial & Monocot & Angiosperm & Hens; Pakanen; J\"ak\"al\"aniemi; Tuomi; Kvist & Biol Conserv & Journal Article & NA & 2017 & 10.1016/j.biocon.2017.04.019 & Hens Biol Conserv 2017 Appendix A. Supplementary Data & 16 & 2000 & 2015 & 1 & No & Yes & No & Ampumavaara; Kiutak\"ong\"as N; Patok\"ong\"as; Mataraniemi; Kiutak\"ong\"as b; Kiutak\"ong\"as E; Kiutak\"ong\"as NWb; Kiutak\"ong\"as NWc; Kiutak\"ong\"as Swa; Kiutak\"ong\"as SWc; Kiutak\"ong\"as W & 12 & NA & NA & NA & FIN & Europe & BOR & A & Mean & No & Unmanipulated & W & 2000 & NA & NA & 2015 & NA & NA & NA & Divided & NA & NA & 8 & 0.989091 & COMPADRE & Released & Released \\ 
238566 & Epipactis\_atrorubens\_3 & Epipactis atrorubens & Darkred helleborine & Plantae & Tracheophyta & Liliopsida & Asparagales & Orchidaceae & Epipactis & atrorubens & NA & NA & Herbaceous perennial & Monocot & Angiosperm & Hens; Pakanen; J\"ak\"al\"aniemi; Tuomi; Kvist & Biol Conserv & Journal Article & NA & 2017 & 10.1016/j.biocon.2017.04.019 & Hens Biol Conserv 2017 Appendix A. Supplementary Data & 16 & 2000 & 2015 & 1 & No & Yes & No & Ampumavaara & 12 & 66.37472 & 29.31944 & 214 & FIN & Europe & BOR & A & Mean & NA & Unmanipulated & W & 2000 & NA & NA & 2015 & NA & NA & NA & Divided & NA & The GPS coordinates were approximated to the closest geographic location described in the reference & 8 & 1.001000 & COMPADRE & Released & Released \\ 
238568 & Epipactis\_atrorubens\_3 & Epipactis atrorubens & Darkred helleborine & Plantae & Tracheophyta & Liliopsida & Asparagales & Orchidaceae & Epipactis & atrorubens & NA & NA & Herbaceous perennial & Monocot & Angiosperm & Hens; Pakanen; J\"ak\"al\"aniemi; Tuomi; Kvist & Biol Conserv & Journal Article & NA & 2017 & 10.1016/j.biocon.2017.04.019 & Hens Biol Conserv 2017 Appendix A. Supplementary Data & 16 & 2000 & 2015 & 1 & No & Yes & No & Patok\"ong\"as & 12 & 66.38944 & 29.26778 & 208 & FIN & Europe & BOR & A & Mean & NA & Unmanipulated & W & 2002 & NA & NA & 2015 & NA & NA & NA & Divided & NA & The GPS coordinates were approximated to the closest geographic location described in the reference & 8 & 1.001000 & COMPADRE & Released & Released \\ 
\bottomrule
\end{tabular*}
\end{table}

\begin{Shaded}
\begin{Highlighting}[]
\CommentTok{\# Nota que en el genero *Epipactis*, tenemos 51 fila, pero todas son de la misma especie de dos estudios distintos.}
\end{Highlighting}
\end{Shaded}

\begin{center}\rule{0.5\linewidth}{0.5pt}\end{center}

\section{Matrices promedio}\label{matrices-promedio}

Para un estudio y especie determinados, la base de datos COMPADRE puede
contener múltiples matrices que reflejan diferentes combinaciones de
población, año y/o tratamiento. Hay una función que permite obtener una
única matriz de promedio general para cada grupo de interés (p.~ej.,
MatrixPopulation) y, por lo tanto, limitar la pseudorreplicación. El
termino usado es colapsar las matrices, con la función mpm\_mean.

Seleccionamos la búsqueda anterior y seleccionamos solamente el estudio
de Hens et al.~en 2017.

El primer paso es convertir cada una de las matrices de transición de
estados de poblaciones en una lista de matrices y luego usar la función
mpm\_mean para calcular la matriz promedio de todas las matrices de
transición de estados de poblaciones. De igual manera se puede calcular
la mediana de los elementos de la matriz con la función mpm\_median.

\begin{Shaded}
\begin{Highlighting}[]
\NormalTok{Epi\_Hens }\OtherTok{=}\NormalTok{ index\_O }\SpecialCharTok{\%\textgreater{}\%}
  \FunctionTok{filter}\NormalTok{(SpeciesAuthor }\SpecialCharTok{\%in\%} \FunctionTok{c}\NormalTok{(}\StringTok{"Epipactis\_atrorubens\_3"}\NormalTok{))   }\CommentTok{\# Filtrar la información de *Epipactis* en la base de datos.}

\NormalTok{mpms }\OtherTok{\textless{}{-}}\NormalTok{ Epi\_Hens}\SpecialCharTok{$}\NormalTok{mat[Epi\_Hens}\SpecialCharTok{$}\NormalTok{SpeciesAuthor }\SpecialCharTok{==} \StringTok{"Epipactis\_atrorubens\_3"}\NormalTok{] }\CommentTok{\# convertir en lista de matrices}

\FunctionTok{mpm\_mean}\NormalTok{(mpms) }\CommentTok{\# calcular la matriz promedio de todas las matrices de transición de estados de poblaciones. Se consigue una matriz para matA, matU, matF y matC.}
\end{Highlighting}
\end{Shaded}

\begin{verbatim}
  MatrixClassOrganized
1               active
2               active
3               active
4               active
5                 dorm
6                 dorm
7                 dorm
8                 dorm
                                                                                                                                MatrixClassAuthor
1                                                                                          S: seedling (smaller than 3 cm; maximum of two leaves)
2                                                                                                  V: vegetative (one vegetative shoot (>= 3 cm))
3 MV: mature vegetative (more than one vegetative shoot; no fertile shoots or a plant with no fertile shoots that flowered in the previous years)
4                                                                                                         F: fertile (at least one fertile shoot)
5                                                                                                                            DS: dormant seedling
6                                                                                                                          DV: dormant vegetative
7                                                                                                                  DMV: dormant mature vegetative
8                                                                                                                             DF: dormant fertile

matA:
          1          2            3          4         5          6          7
1 0.0000000 0.00000000 0.0000000000 0.02707658 0.0000000 0.00000000 0.00000000
2 0.2133690 0.64094625 0.0000000000 0.00000000 0.2121606 0.57071866 0.00000000
3 0.0000000 0.01668078 0.4676567012 0.16649091 0.0000000 0.01921026 0.36824727
4 0.0000000 0.11313765 0.3863594378 0.73172276 0.0000000 0.06576934 0.13218976
5 0.1416013 0.00000000 0.0000000000 0.00000000 0.0995277 0.00000000 0.00000000
6 0.0000000 0.16240300 0.0000000000 0.00000000 0.0000000 0.13653180 0.00000000
7 0.0000000 0.00000000 0.0970043717 0.00000000 0.0000000 0.00000000 0.09879857
8 0.0000000 0.00000000 0.0002894286 0.08886667 0.0000000 0.00000000 0.00000000
             8
1 0.000000e+00
2 0.000000e+00
3 2.671518e-01
4 3.401361e+06
5 0.000000e+00
6 0.000000e+00
7 0.000000e+00
8 1.963378e-01

matU:
          1          2            3          4         5          6           7
1 0.0000000 0.00000000 0.0000000000 0.00000000 0.0000000 0.00000000 0.000000000
2 0.2133690 0.64094625 0.0000000000 0.00000000 0.2121606 0.57071866 0.008163265
3 0.0000000 0.01668078 0.4676567012 0.16649091 0.0000000 0.01921026 0.364165641
4 0.0000000 0.11313765 0.3863594378 0.73172276 0.0000000 0.06576934 0.128108131
5 0.1416013 0.00000000 0.0000000000 0.00000000 0.0995277 0.00000000 0.000000000
6 0.0000000 0.16240300 0.0000000000 0.00000000 0.0000000 0.13653180 0.000000000
7 0.0000000 0.00000000 0.0970043717 0.00000000 0.0000000 0.00000000 0.098798574
8 0.0000000 0.00000000 0.0002894286 0.08886667 0.0000000 0.00000000 0.000000000
            8
1 0.000000000
2 0.005102041
3 0.264600740
4 0.326799803
5 0.000000000
6 0.000000000
7 0.000000000
8 0.196337845

matF:
  1 2 3          4 5 6 7 8
1 0 0 0 0.02707658 0 0 0 0
2 0 0 0 0.00000000 0 0 0 0
3 0 0 0 0.00000000 0 0 0 0
4 0 0 0 0.00000000 0 0 0 0
5 0 0 0 0.00000000 0 0 0 0
6 0 0 0 0.00000000 0 0 0 0
7 0 0 0 0.00000000 0 0 0 0
8 0 0 0 0.00000000 0 0 0 0

matC:
  1 2 3 4 5 6 7 8
1 0 0 0 0 0 0 0 0
2 0 0 0 0 0 0 0 0
3 0 0 0 0 0 0 0 0
4 0 0 0 0 0 0 0 0
5 0 0 0 0 0 0 0 0
6 0 0 0 0 0 0 0 0
7 0 0 0 0 0 0 0 0
8 0 0 0 0 0 0 0 0
\end{verbatim}

\begin{center}\rule{0.5\linewidth}{0.5pt}\end{center}

\section{Calcular la desviación de los elementos de la
matriz}\label{calcular-la-desviaciuxf3n-de-los-elementos-de-la-matriz}

De igual manera se puede calcular la desviación estándar de los
elementos de la matriz de transición de estados de poblaciones con la
función mpm\_sd.

\begin{Shaded}
\begin{Highlighting}[]
\FunctionTok{mpm\_sd}\NormalTok{(mpms) }\CommentTok{\# calcular la desviación estander de los elementos de la matriz de transición de estados de poblaciones. Se consigue una matriz para matA, matU, matF y matC.}
\end{Highlighting}
\end{Shaded}

\begin{verbatim}
  MatrixClassOrganized
1               active
2               active
3               active
4               active
5                 dorm
6                 dorm
7                 dorm
8                 dorm
                                                                                                                                MatrixClassAuthor
1                                                                                          S: seedling (smaller than 3 cm; maximum of two leaves)
2                                                                                                  V: vegetative (one vegetative shoot (>= 3 cm))
3 MV: mature vegetative (more than one vegetative shoot; no fertile shoots or a plant with no fertile shoots that flowered in the previous years)
4                                                                                                         F: fertile (at least one fertile shoot)
5                                                                                                                            DS: dormant seedling
6                                                                                                                          DV: dormant vegetative
7                                                                                                                  DMV: dormant mature vegetative
8                                                                                                                             DF: dormant fertile

matA:
          1          2           3          4         5         6         7
1 0.0000000 0.00000000 0.000000000 0.04774682 0.0000000 0.0000000 0.0000000
2 0.3625333 0.22980249 0.000000000 0.00000000 0.3816246 0.3695790 0.0000000
3 0.0000000 0.02864846 0.302153080 0.14768233 0.0000000 0.0607252 0.4063030
4 0.0000000 0.12468319 0.301021644 0.19016152 0.0000000 0.1488403 0.2284169
5 0.2788961 0.00000000 0.000000000 0.00000000 0.2610421 0.0000000 0.0000000
6 0.0000000 0.16508248 0.000000000 0.00000000 0.0000000 0.2014153 0.0000000
7 0.0000000 0.00000000 0.121154533 0.00000000 0.0000000 0.0000000 0.2191404
8 0.0000000 0.00000000 0.001861297 0.11273170 0.0000000 0.0000000 0.0000000
             8
1 0.000000e+00
2 0.000000e+00
3 3.151303e-01
4 2.380952e+07
5 0.000000e+00
6 0.000000e+00
7 0.000000e+00
8 2.411732e-01

matU:
          1          2           3         4         5         6          7
1 0.0000000 0.00000000 0.000000000 0.0000000 0.0000000 0.0000000 0.00000000
2 0.3625333 0.22980249 0.000000000 0.0000000 0.3816246 0.3695790 0.05714286
3 0.0000000 0.02864846 0.302153080 0.1476823 0.0000000 0.0607252 0.40698137
4 0.0000000 0.12468319 0.301021644 0.1901615 0.0000000 0.1488403 0.22896623
5 0.2788961 0.00000000 0.000000000 0.0000000 0.2610421 0.0000000 0.00000000
6 0.0000000 0.16508248 0.000000000 0.0000000 0.0000000 0.2014153 0.00000000
7 0.0000000 0.00000000 0.121154533 0.0000000 0.0000000 0.0000000 0.21914041
8 0.0000000 0.00000000 0.001861297 0.1127317 0.0000000 0.0000000 0.00000000
           8
1 0.00000000
2 0.03571429
3 0.31577729
4 0.36251239
5 0.00000000
6 0.00000000
7 0.00000000
8 0.24117318

matF:
  1 2 3          4 5 6 7 8
1 0 0 0 0.04774682 0 0 0 0
2 0 0 0 0.00000000 0 0 0 0
3 0 0 0 0.00000000 0 0 0 0
4 0 0 0 0.00000000 0 0 0 0
5 0 0 0 0.00000000 0 0 0 0
6 0 0 0 0.00000000 0 0 0 0
7 0 0 0 0.00000000 0 0 0 0
8 0 0 0 0.00000000 0 0 0 0

matC:
  1 2 3 4 5 6 7 8
1 0 0 0 0 0 0 0 0
2 0 0 0 0 0 0 0 0
3 0 0 0 0 0 0 0 0
4 0 0 0 0 0 0 0 0
5 0 0 0 0 0 0 0 0
6 0 0 0 0 0 0 0 0
7 0 0 0 0 0 0 0 0
8 0 0 0 0 0 0 0 0
\end{verbatim}

Muchas veces es poco tedioso evaluar la base de datos completa de
COMPADRE. Una manera facilitar el acceso a los datos es extraer las
matrices con la función cdb\_flatten. Nota lo que hace esta función es
extraer todas las matrices de transición de estados de poblaciones
(matA, matU, matF, matC) y las etapas de cada matriz y crear una lista
con esa información.

Aquí hay 49 filas que corresponden a las 49 matrices de transición de
estados de poblaciones de orquídeas del estudio de Hens et al.~2017.

\begin{Shaded}
\begin{Highlighting}[]
\NormalTok{Epi\_flat }\OtherTok{=} \FunctionTok{cdb\_flatten}\NormalTok{(Epi\_Hens) }\CommentTok{\# Extraer las matrices de transición de estados de la población}


\NormalTok{Epi\_flat}\SpecialCharTok{$}\NormalTok{matA[}\DecValTok{1}\NormalTok{] }\CommentTok{\# miramos solamente la primera matriz}
\end{Highlighting}
\end{Shaded}

\begin{verbatim}
[1] "[0 0 0 0.018636 0 0 0 0 0.314 0.723818 0 0 0.190727 0.639182 0 0 0 0.010818 0.448 0.145273 0 0.007636 0.255545 0.308818 0 0.115091 0.412 0.775727 0 0.098455 0.063727 0.433636 0.091182 0 0 0 0.082 0 0 0 0 0.137545 0 0 0 0.073 0 0 0 0 0.046364 0 0 0 0.044273 0 0 0 0.001182 0.068091 0 0 0 0.105727]"
\end{verbatim}

\begin{center}\rule{0.5\linewidth}{0.5pt}\end{center}

\subsection{Evaluación de cumplimiento de
condiciones}\label{evaluaciuxf3n-de-cumplimiento-de-condiciones}

Es sumamente importante el cumplimiento de los supuestos de los
análisis. Si no evaluamos críticamente las condiciones de las matrices,
pudiese ser que los resultados no sean biológicamente relevantes. Para
los análisis de PMP es necesario que las matrices cumplan con ciertas
condiciones y dependiendo de la pregunta científica se tiene que evaluar
las condiciones específicas. Una de estas condiciones es la ergodicidad
e irreducibilidad.

\begin{itemize}
\item
  \textbf{check\_ergodic}: Para la definición de ergodicidad.
\item
  \textbf{check\_irreducible}: Para la definición de irreducibilidad.
\item
  \textbf{check\_NA\_A}: Determinar si hay valores faltantes en `matA'
  Los valores faltantes (``NA'') en las matrices impiden que la mayoría
  de los cálculos utilicen esas matrices.
\item
  \textbf{check\_NA\_U}: Determinar si hay valores faltantes en `matU'
  Los valores faltantes (``NA'') en las matrices impiden que la mayoría
  de los cálculos utilicen esas matrices. La matriz de supervivencia es
  la matriz más importante para evaluar el Caladlargo de vida de una
  especie.
\item
  \textbf{check\_NA\_F}: Determinar si hay valores faltantes en `matF'
  Los valores faltantes (``NA'') en las matrices impiden que la mayoría
  de los cálculos utilicen esas matrices. La matriz de fecundidad es la
  matriz más importante para evaluar el crecimiento de la población,
  cualquier especie donde no tiene fecundidad se verá que tiende
  disminiur a través del tiempo.
\item
  \textbf{check\_NA\_C}: Determinar si hay valores faltantes en `matC'.
  Los valores faltantes (``NA'') en las matrices impiden que la mayoría
  de los cálculos utilicen esas matrices.
\item
  \textbf{check\_zero\_U}: Determinar si la `matU' tiene todos ceros
  (incluido `NA'). Las submatrices compuestas enteramente de valores
  cero pueden resultar problemáticas. Puede haber buenas razones
  biológicas para este fenómeno. Por ejemplo, en la población focal
  particular en el año focal particular, realmente no se registró una
  supervivencia. Sin embargo, las submatrices de valor cero pueden
  provocar que algunos cálculos fallen y puede ser necesario excluirlas.
\item
  \textbf{check\_zero\_F}: Determinar si la `matF' tiene todos ceros
  (incluido `NA'). Las submatrices compuestas enteramente de valores
  cero pueden resultar problemáticas. Puede haber buenas razones
  biológicas para este fenómeno o que se asocien con preguntas
  científicas específicas. Por ejemplo, en la población particular en el
  año particular, realmente no se registró reproducción. Sin embargo,
  las submatrices con valores de cero pueden provocar que algunos
  cálculos fallen y puede ser necesario excluirlas. En otros estudios no
  se evaluó la fecundidad ya que estos parámetros no eran necesarios
  para contestar las preguntas científicas.
\item
  \textbf{check\_zero\_U\_colsum}: Determinar si las columnas de `matU'
  suman a cero, esto implica que no hay supervivencia de esa etapa en
  particular. Esta puede ser una parametrización perfectamente válida
  para un año/lugar en particular, pero es biológicamente irracional a
  largo plazo y los usuarios pueden desear excluir matrices
  problemáticas de su análisis.
\item
  \textbf{check\_singular\_U}: Determinar si `matU' es singular. Se dice
  que las matrices son singulares si no se pueden invertir. Se requiere
  inversión para muchos cálculos matriciales y, por lo tanto, la
  singularidad puede hacer que algunos cálculos fallen.
\item
  \textbf{check\_component\_sum}: Determinar si las submatrices
  `matU'/`matF'/`matC' suman a `matA'. Un MPM completo (``matA'') se
  puede dividir en las submatrices que lo componen (es decir, ``matU'',
  ``matF'' y ``matC''). La suma de estas submatrices debe ser igual al
  MPM completo (es decir, `matA' = `matU' + `matF' + `matC'). A veces,
  sin embargo, se producen errores de modo que las submatrices NO suman
  `matA'. Normalmente, esto se debe a errores de redondeo, pero es
  posible que se produzcan errores más importantes.
\item
  \textbf{check\_ergodic}: Determinar si `matA' es ergódico (ver
  isErgodic). Algunos cálculos matriciales requieren que el MPM
  (``matA'') sea ergódico. Los MPM ergódicos son aquellos en los que
  existe un único estado estable asintótico que no depende de la
  estructura de la etapa inicial. Por el contrario, los MPM no ergódicos
  son aquellos en los que existen múltiples estados estables
  asintóticos, que dependen de la estructura de la etapa inicial. Los
  MPM que no son ergódicos suelen ser biológicamente irracionales, tanto
  en términos de la descripción de su ciclo de vida como de su dinámica
  proyectada. Por tanto, hacen que algunos cálculos fallen.
\item
  \textbf{check\_irreductible}: Determinar si `matA' es irreducible (ver
  isIrreducible). Algunos cálculos matriciales requieren que el MPM
  (``matA'') sea irreducible. Los MPM irreductibles son aquellos en los
  que las tasas de transición parametrizadas facilitan el paso de todas
  las etapas a todas las demás etapas. Por el contrario, los MPM
  reducibles representan ciclos de vida incompletos donde no son
  posibles caminos desde todas las etapas a todas las demás etapas. Los
  MPM que son reducibles suelen ser biológicamente irracionales (pero no
  siempre), tanto en términos de la descripción de su ciclo de vida como
  de su dinámica proyectada. Por ello, hacen que algunos cálculos
  fallen. La irreductibilidad es necesaria pero no suficiente para la
  ergodicidad.
\item
  \textbf{check\_primitive}: Deteminar si `matA' es primitivo (ver
  isPrimitive). Una matriz primitiva es una matriz no negativa que es
  irreducible y tiene un solo valor propio de módulo máximo. Por lo
  tanto, esta verificación es redundante debido a la superposición con
  `check\_irreducible' y `checkErdogic'.
\item
  \textbf{check\_surv\_gte\_1}: Determinar si `matU' contiene valores
  iguales o mayores que 1. La supervivencia está limitada entre 0 y 1.
  Los valores superiores a 1 no son biológicamente razonables (vea
  capítulo Impacto de Datos sin sentidos).
\end{itemize}

Usando la función cdb\_flag se puede evaluar todas las opciones en la
base de datos de COMPADRE.

Usando la función \textbf{cdb\_flag} se identifican problemas
potenciales en las matrices de transición de estados de las poblaciones.
Esa función añade una fila por opciones en la matriz. Nota que para cada
prueba anterior crea una columna con \textbf{TRUE} o \textbf{FALSE} para
cada matriz en la base de datos.

Para facilitar el análisis usamos solamente las especies de
\emph{Caladenia} y enseñamos solamente las tres primeras filas.

\begin{Shaded}
\begin{Highlighting}[]
\NormalTok{Caladenia }\OtherTok{=}\NormalTok{ index\_O }\SpecialCharTok{\%\textgreater{}\%}
  \FunctionTok{filter}\NormalTok{(Genus }\SpecialCharTok{\%in\%} \FunctionTok{c}\NormalTok{(}\StringTok{"Caladenia"}\NormalTok{)) }\CommentTok{\# Filtrar la información de *Caladenia* en la base de datos.}



\NormalTok{Compadre\_flagged }\OtherTok{\textless{}{-}} \FunctionTok{cdb\_flag}\NormalTok{(Caladenia) }\CommentTok{\# Evaluación de todas las opciones en la base de datos de COMPADRE}

\NormalTok{flagged}\OtherTok{=}\NormalTok{Compadre\_flagged }\SpecialCharTok{\%\textgreater{}\%} \FunctionTok{select}\NormalTok{(mat, SpeciesAccepted, }\FunctionTok{starts\_with}\NormalTok{(}\StringTok{"check"}\NormalTok{))}

\FunctionTok{head}\NormalTok{(}\FunctionTok{cdb\_metadata}\NormalTok{(flagged), }\AttributeTok{n =} \DecValTok{3}\NormalTok{) }\SpecialCharTok{|\textgreater{}} 
  \FunctionTok{gt}\NormalTok{()}
\end{Highlighting}
\end{Shaded}

\begin{table}
\fontsize{12.0pt}{14.0pt}\selectfont
\begin{tabular*}{\linewidth}{@{\extracolsep{\fill}}lcccccccccccccc}
\toprule
SpeciesAccepted & check\_NA\_A & check\_NA\_U & check\_NA\_F & check\_NA\_C & check\_zero\_U & check\_zero\_F & check\_zero\_C & check\_zero\_U\_colsum & check\_singular\_U & check\_component\_sum & check\_ergodic & check\_irreducible & check\_primitive & check\_surv\_gte\_1 \\ 
\midrule\addlinespace[2.5pt]
Caladenia amonea & FALSE & FALSE & FALSE & FALSE & FALSE & TRUE & TRUE & FALSE & FALSE & TRUE & TRUE & TRUE & TRUE & FALSE \\ 
Caladenia argocalla & FALSE & FALSE & FALSE & FALSE & FALSE & TRUE & TRUE & FALSE & FALSE & TRUE & TRUE & TRUE & TRUE & FALSE \\ 
Caladenia clavigera & FALSE & FALSE & FALSE & FALSE & FALSE & TRUE & TRUE & FALSE & FALSE & TRUE & TRUE & TRUE & TRUE & FALSE \\ 
\bottomrule
\end{tabular*}
\end{table}

\begin{center}\rule{0.5\linewidth}{0.5pt}\end{center}

\subsection{Ergodicidad e
Irreductibilidad}\label{ergodicidad-e-irreductibilidad}

En los siguientes procedimientos se aplican filtros específicos a
matrices de transición poblacional con el objetivo de evaluar dos
propiedades fundamentales en el análisis demográfico: la ergodicidad y
la irreductibilidad. Se observa que toda matriz ergódica cumple con la
condición de ser irreducible; sin embargo, no todas las matrices
irreducibles son necesariamente ergódicas. Para ilustrar esta relación,
se cuantificó el número de matrices que cumplen con cada propiedad,
permitiendo así una comparación entre ambos conjuntos.

\subsubsection{Matrices ergódicas}\label{matrices-erguxf3dicas}

\begin{Shaded}
\begin{Highlighting}[]
\NormalTok{Mat\_ergodic }\OtherTok{\textless{}{-}} \FunctionTok{subset}\NormalTok{(}
\NormalTok{  Compadre\_flagged,}
\NormalTok{  check\_NA\_A }\SpecialCharTok{==} \ConstantTok{FALSE} \SpecialCharTok{\&}\NormalTok{ check\_ergodic }\SpecialCharTok{==} \ConstantTok{TRUE}
\NormalTok{) }\CommentTok{\# para evaluar si las matrices son ergódicas solamente}
\FunctionTok{head}\NormalTok{(}\FunctionTok{cdb\_metadata}\NormalTok{(Mat\_ergodic), }\AttributeTok{n=}\DecValTok{2}\NormalTok{) }\SpecialCharTok{|\textgreater{}} 
  \FunctionTok{gt}\NormalTok{()}
\end{Highlighting}
\end{Shaded}

\begin{table}
\fontsize{12.0pt}{14.0pt}\selectfont
\begin{tabular*}{\linewidth}{@{\extracolsep{\fill}}rlllllllllllllllllllrllrrrrllllrrrrllllllllrlrrlrllllrrlllcccccccccccccc}
\toprule
MatrixID & SpeciesAuthor & SpeciesAccepted & CommonName & Kingdom & Phylum & Class & Order & Family & Genus & Species & Infraspecies & InfraspeciesType & OrganismType & DicotMonoc & AngioGymno & Authors & Journal & SourceType & OtherType & YearPublication & DOI\_ISBN & AdditionalSource & StudyDuration & StudyStart & StudyEnd & ProjectionInterval & MatrixCriteriaSize & MatrixCriteriaOntogeny & MatrixCriteriaAge & MatrixPopulation & NumberPopulations & Lat & Lon & Altitude & Country & Continent & Ecoregion & StudiedSex & MatrixComposite & MatrixSeasonal & MatrixTreatment & MatrixCaptivity & MatrixStartYear & MatrixStartSeason & MatrixStartMonth & MatrixEndYear & MatrixEndSeason & MatrixEndMonth & CensusType & MatrixSplit & MatrixFec & Observations & MatrixDimension & SurvivalIssue & \_Database & \_PopulationStatus & \_PublicationStatus & check\_NA\_A & check\_NA\_U & check\_NA\_F & check\_NA\_C & check\_zero\_U & check\_zero\_F & check\_zero\_C & check\_zero\_U\_colsum & check\_singular\_U & check\_component\_sum & check\_ergodic & check\_irreducible & check\_primitive & check\_surv\_gte\_1 \\ 
\midrule\addlinespace[2.5pt]
238285 & Caladenia\_amonea & Caladenia amonea & NA & Plantae & Tracheophyta & Liliopsida & Asparagales & Orchidaceae & Caladenia & amonea & NA & NA & Herbaceous perennial & Monocot & Angiosperm & Raymond L. Tremblay; P\\'erez; Larcombe; Brown; Quarmby; Bickerton; French; Bould & Aust J Bot & Journal Article & NA & 2009 & 10.1071/BT08167 & NA & 12 & 1996 & 2007 & 1 & No & Yes & No & North-eastern Melbourne, Vic & 1 & NA & NA & NA & AUS & Oceania & TBM & A & Mean & No & Unmanipulated & W & 1996 & Summer & 8 & 2007 & Autumn & 9 & NA & Divided & NA & Locations of surveyed areas are kept vague to reduce the likelihood of unauthorised disturbance of the populations & 3 & 1 & COMPADRE & Released & Released & FALSE & FALSE & FALSE & FALSE & FALSE & TRUE & TRUE & FALSE & FALSE & TRUE & TRUE & TRUE & TRUE & FALSE \\ 
238286 & Caladenia\_argocalla & Caladenia argocalla & NA & Plantae & Tracheophyta & Liliopsida & Asparagales & Orchidaceae & Caladenia & argocalla & NA & NA & Herbaceous perennial & Monocot & Angiosperm & Raymond L. Tremblay; P\\'erez; Larcombe; Brown; Quarmby; Bickerton; French; Bould & Aust J Bot & Journal Article & NA & 2009 & 10.1071/BT08167 & NA & 12 & 1996 & 2007 & 1 & No & Yes & No & Adelaide Hills and Clare Valley, SA & 1 & NA & NA & NA & AUS & Oceania & TBM & A & Mean & No & Unmanipulated & W & 2003 & Autumn & 9 & 2007 & Autumn & 10 & NA & Divided & No & Locations of surveyed areas are kept vague to reduce the likelihood of unauthorised disturbance of the populations & 3 & 1 & COMPADRE & Released & Released & FALSE & FALSE & FALSE & FALSE & FALSE & TRUE & TRUE & FALSE & FALSE & TRUE & TRUE & TRUE & TRUE & FALSE \\ 
\bottomrule
\end{tabular*}
\end{table}

\subsubsection{Matrices irreductibles}\label{matrices-irreductibles}

\begin{Shaded}
\begin{Highlighting}[]
\NormalTok{Mat\_irred }\OtherTok{\textless{}{-}} \FunctionTok{subset}\NormalTok{(}
\NormalTok{  Compadre\_flagged,}
\NormalTok{  check\_NA\_A }\SpecialCharTok{==} \ConstantTok{FALSE} \SpecialCharTok{\&}\NormalTok{ check\_irreducible }\SpecialCharTok{==} \ConstantTok{TRUE}
\NormalTok{) }\CommentTok{\# para evaluar si las matrices son irreducibles solamente}

\FunctionTok{head}\NormalTok{(}\FunctionTok{cdb\_metadata}\NormalTok{(Mat\_irred), }\AttributeTok{n=}\DecValTok{2}\NormalTok{) }\SpecialCharTok{|\textgreater{}} 
  \FunctionTok{gt}\NormalTok{()}
\end{Highlighting}
\end{Shaded}

\begin{table}
\fontsize{12.0pt}{14.0pt}\selectfont
\begin{tabular*}{\linewidth}{@{\extracolsep{\fill}}rlllllllllllllllllllrllrrrrllllrrrrllllllllrlrrlrllllrrlllcccccccccccccc}
\toprule
MatrixID & SpeciesAuthor & SpeciesAccepted & CommonName & Kingdom & Phylum & Class & Order & Family & Genus & Species & Infraspecies & InfraspeciesType & OrganismType & DicotMonoc & AngioGymno & Authors & Journal & SourceType & OtherType & YearPublication & DOI\_ISBN & AdditionalSource & StudyDuration & StudyStart & StudyEnd & ProjectionInterval & MatrixCriteriaSize & MatrixCriteriaOntogeny & MatrixCriteriaAge & MatrixPopulation & NumberPopulations & Lat & Lon & Altitude & Country & Continent & Ecoregion & StudiedSex & MatrixComposite & MatrixSeasonal & MatrixTreatment & MatrixCaptivity & MatrixStartYear & MatrixStartSeason & MatrixStartMonth & MatrixEndYear & MatrixEndSeason & MatrixEndMonth & CensusType & MatrixSplit & MatrixFec & Observations & MatrixDimension & SurvivalIssue & \_Database & \_PopulationStatus & \_PublicationStatus & check\_NA\_A & check\_NA\_U & check\_NA\_F & check\_NA\_C & check\_zero\_U & check\_zero\_F & check\_zero\_C & check\_zero\_U\_colsum & check\_singular\_U & check\_component\_sum & check\_ergodic & check\_irreducible & check\_primitive & check\_surv\_gte\_1 \\ 
\midrule\addlinespace[2.5pt]
238285 & Caladenia\_amonea & Caladenia amonea & NA & Plantae & Tracheophyta & Liliopsida & Asparagales & Orchidaceae & Caladenia & amonea & NA & NA & Herbaceous perennial & Monocot & Angiosperm & Raymond L. Tremblay; P\\'erez; Larcombe; Brown; Quarmby; Bickerton; French; Bould & Aust J Bot & Journal Article & NA & 2009 & 10.1071/BT08167 & NA & 12 & 1996 & 2007 & 1 & No & Yes & No & North-eastern Melbourne, Vic & 1 & NA & NA & NA & AUS & Oceania & TBM & A & Mean & No & Unmanipulated & W & 1996 & Summer & 8 & 2007 & Autumn & 9 & NA & Divided & NA & Locations of surveyed areas are kept vague to reduce the likelihood of unauthorised disturbance of the populations & 3 & 1 & COMPADRE & Released & Released & FALSE & FALSE & FALSE & FALSE & FALSE & TRUE & TRUE & FALSE & FALSE & TRUE & TRUE & TRUE & TRUE & FALSE \\ 
238286 & Caladenia\_argocalla & Caladenia argocalla & NA & Plantae & Tracheophyta & Liliopsida & Asparagales & Orchidaceae & Caladenia & argocalla & NA & NA & Herbaceous perennial & Monocot & Angiosperm & Raymond L. Tremblay; P\\'erez; Larcombe; Brown; Quarmby; Bickerton; French; Bould & Aust J Bot & Journal Article & NA & 2009 & 10.1071/BT08167 & NA & 12 & 1996 & 2007 & 1 & No & Yes & No & Adelaide Hills and Clare Valley, SA & 1 & NA & NA & NA & AUS & Oceania & TBM & A & Mean & No & Unmanipulated & W & 2003 & Autumn & 9 & 2007 & Autumn & 10 & NA & Divided & No & Locations of surveyed areas are kept vague to reduce the likelihood of unauthorised disturbance of the populations & 3 & 1 & COMPADRE & Released & Released & FALSE & FALSE & FALSE & FALSE & FALSE & TRUE & TRUE & FALSE & FALSE & TRUE & TRUE & TRUE & TRUE & FALSE \\ 
\bottomrule
\end{tabular*}
\end{table}

\subsubsection{Matrices ergódicas e
irreductibles}\label{matrices-erguxf3dicas-e-irreductibles}

\begin{Shaded}
\begin{Highlighting}[]
\NormalTok{Mat\_erg\_irred }\OtherTok{\textless{}{-}} \FunctionTok{subset}\NormalTok{(}
\NormalTok{  Compadre\_flagged,}
\NormalTok{  check\_NA\_A }\SpecialCharTok{==} \ConstantTok{FALSE} \SpecialCharTok{\&}\NormalTok{ check\_irreducible }\SpecialCharTok{==} \ConstantTok{TRUE} \SpecialCharTok{\&}\NormalTok{ check\_ergodic }\SpecialCharTok{==} \ConstantTok{TRUE}
\NormalTok{) }\CommentTok{\# para evaluar si las matrices son ergódicas e irreducibles al mismo tiempo}

\FunctionTok{head}\NormalTok{(}\FunctionTok{cdb\_metadata}\NormalTok{(Mat\_erg\_irred), }\AttributeTok{n=}\DecValTok{2}\NormalTok{) }\SpecialCharTok{|\textgreater{}} 
  \FunctionTok{gt}\NormalTok{()}
\end{Highlighting}
\end{Shaded}

\begin{table}
\fontsize{12.0pt}{14.0pt}\selectfont
\begin{tabular*}{\linewidth}{@{\extracolsep{\fill}}rlllllllllllllllllllrllrrrrllllrrrrllllllllrlrrlrllllrrlllcccccccccccccc}
\toprule
MatrixID & SpeciesAuthor & SpeciesAccepted & CommonName & Kingdom & Phylum & Class & Order & Family & Genus & Species & Infraspecies & InfraspeciesType & OrganismType & DicotMonoc & AngioGymno & Authors & Journal & SourceType & OtherType & YearPublication & DOI\_ISBN & AdditionalSource & StudyDuration & StudyStart & StudyEnd & ProjectionInterval & MatrixCriteriaSize & MatrixCriteriaOntogeny & MatrixCriteriaAge & MatrixPopulation & NumberPopulations & Lat & Lon & Altitude & Country & Continent & Ecoregion & StudiedSex & MatrixComposite & MatrixSeasonal & MatrixTreatment & MatrixCaptivity & MatrixStartYear & MatrixStartSeason & MatrixStartMonth & MatrixEndYear & MatrixEndSeason & MatrixEndMonth & CensusType & MatrixSplit & MatrixFec & Observations & MatrixDimension & SurvivalIssue & \_Database & \_PopulationStatus & \_PublicationStatus & check\_NA\_A & check\_NA\_U & check\_NA\_F & check\_NA\_C & check\_zero\_U & check\_zero\_F & check\_zero\_C & check\_zero\_U\_colsum & check\_singular\_U & check\_component\_sum & check\_ergodic & check\_irreducible & check\_primitive & check\_surv\_gte\_1 \\ 
\midrule\addlinespace[2.5pt]
238285 & Caladenia\_amonea & Caladenia amonea & NA & Plantae & Tracheophyta & Liliopsida & Asparagales & Orchidaceae & Caladenia & amonea & NA & NA & Herbaceous perennial & Monocot & Angiosperm & Raymond L. Tremblay; P\\'erez; Larcombe; Brown; Quarmby; Bickerton; French; Bould & Aust J Bot & Journal Article & NA & 2009 & 10.1071/BT08167 & NA & 12 & 1996 & 2007 & 1 & No & Yes & No & North-eastern Melbourne, Vic & 1 & NA & NA & NA & AUS & Oceania & TBM & A & Mean & No & Unmanipulated & W & 1996 & Summer & 8 & 2007 & Autumn & 9 & NA & Divided & NA & Locations of surveyed areas are kept vague to reduce the likelihood of unauthorised disturbance of the populations & 3 & 1 & COMPADRE & Released & Released & FALSE & FALSE & FALSE & FALSE & FALSE & TRUE & TRUE & FALSE & FALSE & TRUE & TRUE & TRUE & TRUE & FALSE \\ 
238286 & Caladenia\_argocalla & Caladenia argocalla & NA & Plantae & Tracheophyta & Liliopsida & Asparagales & Orchidaceae & Caladenia & argocalla & NA & NA & Herbaceous perennial & Monocot & Angiosperm & Raymond L. Tremblay; P\\'erez; Larcombe; Brown; Quarmby; Bickerton; French; Bould & Aust J Bot & Journal Article & NA & 2009 & 10.1071/BT08167 & NA & 12 & 1996 & 2007 & 1 & No & Yes & No & Adelaide Hills and Clare Valley, SA & 1 & NA & NA & NA & AUS & Oceania & TBM & A & Mean & No & Unmanipulated & W & 2003 & Autumn & 9 & 2007 & Autumn & 10 & NA & Divided & No & Locations of surveyed areas are kept vague to reduce the likelihood of unauthorised disturbance of the populations & 3 & 1 & COMPADRE & Released & Released & FALSE & FALSE & FALSE & FALSE & FALSE & TRUE & TRUE & FALSE & FALSE & TRUE & TRUE & TRUE & TRUE & FALSE \\ 
\bottomrule
\end{tabular*}
\end{table}

http://127.0.0.1:29465/rmd\_output/1/compadre\_orchid-ppm-y-su-uso.html

Revisión RLT:

Sept 15, 2024

Jan 20, 2025

Sept 9, 2025


\backmatter


\end{document}
